\documentclass[11pt,a4paper]{article}
\usepackage[T1]{fontenc}
\usepackage[utf8]{inputenc}
\usepackage{amsmath,amssymb,amsthm}
\usepackage[margin=1in]{geometry}
\usepackage[colorlinks=true,linkcolor=blue,urlcolor=blue]{hyperref}
\theoremstyle{plain}
\newtheorem{theorem}{Theorem}
\newtheorem{lemma}[theorem]{Lemma}
\newtheorem{proposition}[theorem]{Proposition}
\newtheorem{corollary}[theorem]{Corollary}
\theoremstyle{definition}
\newtheorem{remark}[theorem]{Remark}

\title{The empirical recurrence for OEIS A183685, proved by transfer matrix}
\author{Adrian Perez Fontelles\\ \small Independent researcher}
\date{2 September 2026}
\begin{document}
\maketitle

\begin{abstract}
OEIS A183685 counts arrays over $\{0,\dots,1\}$ under a condition on the determinant of every
$2\times2$ subblock, and carries an empirical recurrence of order $8$ contributed by
R.~H.~Hardin. It is true, and it is decidable rather than empirical. The condition is local to
a bounded window of consecutive rows, so the count is a walk count on $S=128$ states and is
$C$-finite; whether the conjectured recurrence annihilates it is then settled by a finite exact
computation, with the Cayley--Hamilton theorem bounding the work.
\end{abstract}

\noindent\small 2020 Mathematics Subject Classification. 05A15, 05B45, 15B36.\normalsize

\section{The conjecture}

OEIS A183685 is ``Number of (n+1) X 7 binary arrays with every 2 X 2 subblock nonsingular.''. It has offset $1$ and begins
\[
192, 1360, 8704, 58448, 382848, 2542368, 16748160, 110871040,\ \dots
\]
The entry states, as a conjecture contributed by its author and never marked settled:
\begin{quote}
Empirical: a(n) = 3*a(n-1) + 48*a(n-2) - 78*a(n-3) - 720*a(n-4) + 612*a(n-5) + 4232*a(n-6) - 1392*a(n-7) - 8192*a(n-8).
\end{quote}
As of the ``Last modified'' line on the live entry (Apr 04 2018 10:48:18, revision 9) this is still
recorded as empirical, and nothing on the entry records it as proved.

\section{What is being counted}

Write $\sigma$ for the determinant of a $2\times2$ subblock
$\begin{pmatrix}p&q\\r&s\end{pmatrix}$, that is $\sigma=ps-qr$. The subblocks of an
$(n+1)\times7$ array over $\{0,\dots,1\}$ form a grid of $n$ rows and $6$
columns, and $\sigma$ colours that grid.

The condition is on each subblock separately: every $\sigma$ must not vanish, which is what
``nonsingular'' says of a $2\times2$ matrix. Nothing
relates one subblock to another.

\section{The count is a walk count}

Because the condition looks at one subblock at a time and a subblock is
built from two consecutive array rows, a single row is a state: rows $r$ and $s$ may follow one
another exactly when all $6$ of the subblocks they form pass. An $(n+1)$-row array is then a
walk of $n$ steps and
\[
a(n)\;=\;\iota^{\!\top}M^{\,n}\tau,\qquad\iota=\tau=(1,\dots,1)^{\!\top},
\]
on the $S=128$ rows.

At $n=1$ the array has two rows and the grid is a single row of $6$
subblocks, each tested on its own, so $a(1)$ counts the
$2\times7$ arrays over $\{0,\dots,1\}$ all of whose subblocks pass. The entry's
first term is $192$, and the model reproduces it along with every later term; that is the
cheapest check that the grid has been read with the right shape.

Over the alphabet $\{0,\dots,1\}$ the determinant of a $2\times2$ subblock takes values in
$[-1,1]$, so the colouring of the grid it induces has at most $3$ values.
That bound is what keeps the construction below finite in the only place where the value of
$\sigma$, rather than a comparison between two of them, has to be carried.
In particular $a$ is $C$-finite. The alignment of the walk index with the entry's own $n$ was
checked against its published terms rather than assumed.

\section{The proof}

\begin{lemma}\label{lem:crit}
Let $q(t)=t^{8}-\sum_i c_it^{8-i}$ be the characteristic polynomial of the
conjectured recurrence. The residual $a(n)-\sum_ic_ia(n-i)$ equals
$\iota^{\!\top}M^{\,j}q(M)\tau$ for the corresponding walk index $j$, so the recurrence
holds from some point on if and only if $\iota^{\!\top}M^{j}q(M)\tau=0$ for all large $j$,
and it is enough to examine $j<S$.
\end{lemma}

\begin{proof}
Factor $M^{j}$ out of $M^{j+8}-\sum_ic_iM^{j+8-i}$. For the bound, $M^{S}$
is an integer combination of $I,M,\dots,M^{S-1}$ by Cayley--Hamilton, so the numbers
$u_j=\iota^{\!\top}M^{j}q(M)\tau$ obey the monic recurrence given by the characteristic
polynomial of $M$; $S$ consecutive zeros force every later one, and the last nonzero $u_j$
pins the threshold exactly.
\end{proof}

\begin{theorem}
\[
a(n)\;=\;3\,a(n-1) + 48\,a(n-2) - 78\,a(n-3) - 720\,a(n-4) + 612\,a(n-5) + 4232\,a(n-6) - 1392\,a(n-7) - 8192\,a(n-8)
\]
for every $n>8$.
\end{theorem}

\begin{proof}
The vector $w=q(M)\tau$ was formed by $8$ matrix--vector products in exact integer
arithmetic and $\iota^{\!\top}M^{j}w$ evaluated for $j=0,1,\dots$, again exactly; those
integers vanish from the index corresponding to $n=8+1$ onwards and the run continued
until the working vector was identically zero, which settles every larger $j$ at once.
\end{proof}

\section{Verification}

Three checks, each independent of the algebra above.

First, the model reproduces all $18$ terms the entry publishes, exactly, in integer
arithmetic. This is what ties the model to the entry: the digraph is built by reading the
entry's English, and a misreading gives different counts.

Second, the conjectured recurrence was evaluated directly on those published terms, with no
matrices involved, and holds wherever the proved range and the data overlap.

Third, the annihilation test of Lemma~\ref{lem:crit} was carried out in exact integer
arithmetic on the whole vector rather than on a sampled prefix.

\begin{thebibliography}{9}
\bibitem{oeis} The OEIS Foundation, \emph{The On-Line Encyclopedia of Integer
Sequences}, \url{https://oeis.org}, 2026. Sequence A183685.
\bibitem{stanley} R.~P.~Stanley, \emph{Enumerative Combinatorics}, Volume 1, 2nd ed.,
Cambridge University Press, 2011. (Transfer-matrix method, Section 4.7.)
\bibitem{minc} H.~Minc, \emph{Permanents}, Addison--Wesley, 1978.
\bibitem{horn} R.~A.~Horn and C.~R.~Johnson, \emph{Matrix Analysis}, 2nd ed., Cambridge
University Press, 2013.
\end{thebibliography}

\end{document}
